% !TEX root = ../Report.tex
\section{Introduction}
\label{sec:introduction}

Drug--target interaction (DTI) prediction aims to identify likely associations between drug molecules and biological targets such as proteins, enzymes, and receptors. Reliable DTI evidence supports drug repurposing, therapeutic mechanism analysis, side-effect investigation, and candidate prioritization. Experimental screening remains essential, but it is expensive and cannot exhaustively evaluate the large space of possible drug--target pairs. Computational models therefore provide a practical way to rank candidate interactions before laboratory validation.

Classical DTI methods often represent the task as supervised learning over drug and protein descriptors. The benchmark study of Yamanishi et al. framed the problem by integrating chemical and genomic spaces \cite{yamanishi2008prediction}, and this view remains influential for descriptor-based prediction. However, biomedical entities are also connected through heterogeneous relations involving diseases, pathways, biological processes, side effects, protein families, and protein--protein interactions. These relations form a biomedical knowledge graph (KG), whose structure can provide evidence that is not present in isolated molecular or sequence features.

This paper studies a hybrid KGE-NFM pipeline for DTI prediction. The first stage learns drug and protein embeddings from a biomedical KG using a knowledge graph embedding (KGE) model. The second stage combines these embeddings with Morgan fingerprints for drugs and CTD descriptors for proteins, then predicts interaction labels with a neural factorization machine (NFM). This design treats DTI prediction neither as descriptor classification alone nor as pure graph completion; instead, it lets the downstream predictor model interactions among graph-derived, molecular, and sequence-derived feature fields.

The paper evaluates two graph encoders under the same downstream NFM protocol. DistMult provides an efficient bilinear KGE baseline \cite{yang2015embedding}. CompGCN provides a relation-aware message-passing encoder that composes entity and relation representations before aggregation \cite{vashishth2020composition}. The comparison asks whether a more expressive KG encoder improves downstream DTI prediction when the descriptor fields and classifier are kept fixed.

The main contributions are:
\begin{itemize}
    \item a compact ACM-style formulation of a reproducible KGE-NFM pipeline for DTI prediction;
    \item a controlled comparison between DistMult-NFM and CompGCN-NFM on Yamanishi08 warm-start settings;
    \item an empirical analysis of descriptor baselines, KGE-enhanced models, hyperparameter tuning, and ablation results under ROC-AUC and PR-AUC.
\end{itemize}
